\chapter{Discussion}
\label{chapterlabel5}

This chapter interprets the empirical findings presented in Chapter 4 in relation to the research questions and the wider literature on night-time mobility. It first considers how the identified night-time usage types are organised across Rail and Bus, before examining their relationships with the London Night Workers Classification and wider socio-spatial context. It then discusses selected transport settings that warrant closer investigation, followed by the implications and limitations of the analysis.

\section{Overall Organisation of Night-time Public Transport Use}

London's night-time public transport use is organised through recurrent and geographically patterned rhythms rather than a single, uniformly declining off-peak condition. Within the common 18:00--05:00 analytical window, locations differ in when activity is concentrated, how directional flows are balanced, how far activity continues into the night, how profiles vary across day types and where these patterns occur in the city. These differences are expressed in different ways across Rail and Bus. The result gives empirical substance to accounts of the urban night as a period with its own uneven temporal and spatial organisation rather than simply a lower-intensity continuation of daytime activity \citep{Schwanen2012,VanLiempt2015}.

Taken together, the results show that London's night-time public transport use has clear but mode-specific temporal and spatial organisation. Rail and Bus differ in the dimensions through which this organisation is expressed, while the identified transport types also show systematic but incomplete associations with night-work geography and wider socio-spatial context. The strongest correspondence occurs at the central and peripheral ends of the typologies, whereas intermediate and late-persistent types indicate that transport-use patterns retain temporal and network characteristics not captured by area context alone.

\section{Different Forms of Night-time Organisation in Rail and Bus}

The two modes express night-time organisation through different dominant dimensions. Rail forms a relatively differentiated station typology, whereas Bus varies more continuously along an activity--persistence gradient.

For Rail, the contrast between central departure-oriented and peripheral arrival-oriented stations suggests that London's centre--periphery structure remains visible after conventional commuting hours, re-expressed through evening dispersal and return flows. This broadly accords with previous station typologies identifying central, residential and other functional roles \citep{Gan2020,Verma2021,Zhang2021}. The dedicated night-time window, however, reveals further differentiation. C2 combines relatively balanced entry and exit activity with several major interchange stations, suggesting that network role may also shape directional patterns. Activity intensity and duration likewise diverge: C1 records the highest overall activity but a below-average post-23:00 share, whereas C3 has lower volume but the strongest late-night persistence. These roles also vary across day types, particularly through directional shifts and later continuation on Friday and Saturday nights. This adds a night-time perspective to the intra-week variation identified by \citet{Cheng2024}, showing how such variation is expressed through direction and persistence within the dedicated night-time window. Rail night-time use is therefore organised through spatial position, direction, network role, duration and day type rather than centrality alone.

Bus shows a more continuous structure. Its four profiles share a similar evening decline and differ mainly in activity level and persistence, producing a broad transition from more persistent central and inner areas to earlier-declining outer areas. C2 and C3 occupy overlapping intermediate positions, while weekend variation does not alter this overall structure substantially. Bus is therefore better interpreted as a graded activity--persistence pattern than as four equally distinct functional area types. This is consistent with the more distributed nature of Bus activity, where multiple stops and local functions are combined within each LSOA.

Taken together, Rail concentrates night-time variation into differentiated nodes, while Bus distributes it across a more continuous area-level gradient. Differences in analytical representation may reinforce this contrast, but the two modes nevertheless reveal different dominant forms of night-time organisation.

\section{Systematic but Incomplete Correspondence with Night-work and Wider Area Context}

RQ2 asks how far the transport types are associated with night-work and wider socio-spatial context. The answer is neither that context fully explains the typologies nor that the associations are incidental. Both modes show systematic correspondence with the independently derived London Night Workers Classification \citep{Mavrogeni2025}, but the form of that correspondence differs. Rail clusters are differentiated in their LNWC compositions, whereas Bus clusters show a categorical centre-to-periphery association. Within both modes, the clearest correspondence occurs at the typological endpoints. Central, intensive night-work environments are over-represented around Rail C1 and Bus C1, while peripheral LNWC types are over-represented around Rail C0 and Bus C0. The weaker correspondence among intermediate types therefore shows that night-work geography helps interpret the transport patterns without reproducing them.

The wider contextual indicators clarify what distinguishes these central and peripheral areas. The central types combine intensive observed activity with younger-adult populations, private renting, low household car access and high facility intensity. These areas also combine employment, residential, interchange and amenity-related functions. This extends previous links between station rhythms and surrounding urban functions \citep{Gan2020} by showing that the central night-time types are also embedded in a distinctive residential and mobility context.

The outer types show the opposite contextual pattern: lower observed activity, earlier contraction, more households with dependent children, greater car access and lower facility intensity. Yet these areas also contain above-average resident employment in several sectors associated with night-time work.

Because these measures describe residents' industries of employment rather than workplace locations, they do not show that these workers used the observed services. Instead, the coexistence of weak realised transport use and relatively high resident employment in night-related sectors points to a possible separation between the residential geography of night work and the geography of night-time mobility. This provides a concrete transport-based example of the broader distinction between residence, employment and mobility conditions discussed in previous London research \citep{McArthur2019,Smeds2020}.

Area context is therefore clearly associated with the transport types, but it does not fully explain them. Centrality and night-work intensity account for much of the broad contrast, particularly at the typological endpoints, but not for every temporal role. The incomplete correspondence is itself informative: LNWC captures an important part of the geography of night-time transport, while the transport profiles additionally distinguish mode-specific features such as direction, interchange role and late-night persistence.

\section{Exceptional Types and Night-time Environments Requiring Closer Attention}

Rail C3 provides the clearest example of where transport rhythms cannot be explained by central night-work intensity alone. It has the strongest late-night persistence but is not concentrated in the most facility-intensive central night-work setting. Its inner and eastern catchments instead combine a younger-adult profile with higher health deprivation, unemployment, social renting and accommodation-and-food employment. This combination complicates an account of late-night mobility centred mainly on entertainment districts and high-volume central activity. C3 therefore provides a transport-based example of the broader and more heterogeneous urban night described by \citet{Shaw2014,Shaw2022}, in which work, care, maintenance and local services may follow temporal and spatial geographies that do not coincide with those of nightlife.

One plausible reading is that persistent Rail activity partly reflects on-site or less temporally flexible work, an interpretation compatible with the higher accommodation-and-food employment observed around C3. This remains only one possible explanation: interchange, leisure, mixed land use and late-running services may also contribute to the observed persistence. C3 matters because several observations converge on an environment distinct from both the central high-volume and outer early-declining types, rather than because it establishes a particular passenger purpose.

The outer early-declining Rail and Bus settings deserve attention for a different reason. Weak realised use may coexist with residential concentrations of employment in sectors commonly associated with non-standard hours. Without origin-destination, service-frequency and time-dependent accessibility data, this cannot be labelled unmet demand or inadequate provision. It does, however, define a focused research question: whether early contraction reflects lower underlying travel needs, different modal choices, existing service conditions, dispersed destinations, or some combination of these. These cases are therefore useful for identifying where additional research is needed.

\section{Implications for Night-time Transport Research and Planning}

The findings suggest that night-time public transport should be understood as a temporally differentiated system rather than as a single off-peak condition, consistent with recent work emphasising the temporal heterogeneity of urban accessibility and mobility patterns \citep{Ryan2023}.

The contrast between Rail and Bus further supports a mode-sensitive approach. Rail is organised through differentiated node roles involving direction, station function and persistence, whereas Bus follows a more continuous area-level activity--persistence gradient. Applying a common interpretation across modes would therefore obscure important differences in how night-time transport use is spatially and temporally organised.

More broadly, neither total activity nor spatial centrality alone is sufficient to identify the night-time transport role of a location. In Rail, high-volume central stations are not necessarily the most persistent, showing that intensity and late-night continuation represent distinct planning-relevant dimensions. Together, these findings provide an adaptable framework for examining night-time transport beyond conventional commute-centred and all-day ridership analyses, while also providing a more differentiated basis for planning. By distinguishing temporal and mode-specific transport roles, the typologies can help transport authorities target questions of service continuity, interchange and time-dependent accessibility more precisely across the network.

\section{Limitations and Future Research}

Several limitations define the scope of these interpretations. NUMBAT and BUSTO represent typical day-type profiles rather than continuous observations, so seasonal variation, major events and short-term disruption are not captured. Rail and Bus also differ in spatial unit, temporal resolution and directional measurement, while station catchments and LSOA aggregation simplify the relationship between transport activity and surrounding urban context. Cross-mode comparisons should therefore be interpreted as broad organisational patterns rather than formally equivalent behaviours. The resulting typologies also simplify continuous variation, particularly for Bus where intermediate profiles partly overlap.

A more fundamental limitation is that observed public transport use cannot be equated with underlying transport demand. Smart-card and operational records capture journeys realised under existing service, accessibility and mode-choice conditions. Low observed activity may therefore reflect alternative modes, service constraints or unrealised travel rather than low demand itself. Contextual associations are similarly ecological rather than individual-level or causal: they describe the environments surrounding transport activity but cannot identify who is travelling or why.

Future research could combine these typologies with longer time series, origin--destination or passenger-level data, and measures of service supply and time-dependent accessibility. This would allow analysis to move beyond describing realised night-time mobility towards explaining the mechanisms that produce it and the extent to which observed use reflects underlying demand.
